\documentclass[10pt,a4paper]{amsart}
\usepackage{amsmath,amssymb,amsthm,graphicx}
\usepackage{verbatim}
\newenvironment{reference}{\comment}{\endcomment}
%\newenvironment{reference}{}{}
\newenvironment{slogan}{\comment}{\endcomment}
\newenvironment{history}{\comment}{\endcomment}

% For commutative diagrams we use Xy-pic
\usepackage[all]{xy}

% We use 2cell for 2-commutative diagrams.
\xyoption{2cell}
% Theorem environments.
%
\theoremstyle{plain}
\newtheorem{theorem}[subsection]{Theorem}
\newtheorem{proposition}[subsection]{Proposition}
\newtheorem{lemma}[subsection]{Lemma}

\theoremstyle{definition}
\newtheorem{definition}[subsection]{Definition}
\newtheorem{example}[subsection]{Example}
\newtheorem{exercise}[subsection]{Exercise}
\newtheorem{situation}[subsection]{Situation}

\theoremstyle{remark}
\newtheorem{remark}[subsection]{Remark}
\newtheorem{remarks}[subsection]{Remarks}

\numberwithin{equation}{subsection}

% Macros
%
\def\lim{\mathop{\mathrm{lim}}\nolimits}
\def\colim{\mathop{\mathrm{colim}}\nolimits}
\def\Spec{\mathop{\mathrm{Spec}}}
\def\Hom{\mathop{\mathrm{Hom}}\nolimits}
\def\Ext{\mathop{\mathrm{Ext}}\nolimits}
\def\SheafHom{\mathop{\mathcal{H}\!\mathit{om}}\nolimits}
\def\SheafExt{\mathop{\mathcal{E}\!\mathit{xt}}\nolimits}
\def\Sch{\mathit{Sch}}
\def\Mor{\mathop{\mathrm{Mor}}\nolimits}
\def\Ob{\mathop{\mathrm{Ob}}\nolimits}
\def\Sh{\mathop{\mathit{Sh}}\nolimits}
\def\NL{\mathop{N\!L}\nolimits}
\def\CH{\mathop{\mathrm{CH}}\nolimits}
\def\proetale{{pro\text{-}\acute{e}tale}}
\def\etale{{\acute{e}tale}}
\def\QCoh{\mathit{QCoh}}
\def\Ker{\mathop{\mathrm{Ker}}}
\def\Im{\mathop{\mathrm{Im}}}
\def\Coker{\mathop{\mathrm{Coker}}}
\def\Coim{\mathop{\mathrm{Coim}}}

% Boxtimes
%
\DeclareMathSymbol{\boxtimes}{\mathbin}{AMSa}{"02}

%
% Macros for moduli stacks/spaces
%
\def\QCohstack{\mathcal{QC}\!\mathit{oh}}
\def\Cohstack{\mathcal{C}\!\mathit{oh}}
\def\Spacesstack{\mathcal{S}\!\mathit{paces}}
\def\Quotfunctor{\mathrm{Quot}}
\def\Hilbfunctor{\mathrm{Hilb}}
\def\Curvesstack{\mathcal{C}\!\mathit{urves}}
\def\Polarizedstack{\mathcal{P}\!\mathit{olarized}}
\def\Complexesstack{\mathcal{C}\!\mathit{omplexes}}
% \Pic is the operator that assigns to X its picard group, usage \Pic(X)
% \Picardstack_{X/B} denotes the Picard stack of X over B
% \Picardfunctor_{X/B} denotes the Picard functor of X over B
\def\Pic{\mathop{\mathrm{Pic}}\nolimits}
\def\Picardstack{\mathcal{P}\!\mathit{ic}}
\def\Picardfunctor{\mathrm{Pic}}
\def\Deformationcategory{\mathcal{D}\!\mathit{ef}}
\usepackage{cleveref}
\crefname{originalreference}{original reference}{original references}
\Crefname{originalreference}{Original reference}{Original references}
\makeatletter
\newlabel{situation-shriek}{{0F42}{source}{}{sourcepointer-1}{}}
\newlabel{situation-shriek@cref}{{[originalreference][0][]0F42}{[1][0][]source}{}{}{}}
\makeatother
\crefname{originallemma}{original lemma}{original lemmas}
\Crefname{originallemma}{Original lemma}{Original lemmas}
\makeatletter
\newlabel{lemma-shriek-affine-line}{{0AA1}{source}{}{sourcepointer-2}{}}
\newlabel{lemma-shriek-affine-line@cref}{{[originallemma][0][]0AA1}{[1][0][]source}{}{}{}}
\makeatother
\crefname{originallemma}{original lemma}{original lemmas}
\Crefname{originallemma}{Original lemma}{Original lemmas}
\makeatletter
\newlabel{dualizing-lemma-dualizing-polynomial-ring}{{0A7J}{source}{}{sourcepointer-3}{}}
\newlabel{dualizing-lemma-dualizing-polynomial-ring@cref}{{[originallemma][0][]0A7J}{[1][0][]source}{}{}{}}
\makeatother
\crefname{originallemma}{original lemma}{original lemmas}
\Crefname{originallemma}{Original lemma}{Original lemmas}
\makeatletter
\newlabel{dualizing-lemma-dualizing-quotient}{{0A7I}{source}{}{sourcepointer-4}{}}
\newlabel{dualizing-lemma-dualizing-quotient@cref}{{[originallemma][0][]0A7I}{[1][0][]source}{}{}{}}
\makeatother
\crefname{originallemma}{original lemma}{original lemmas}
\Crefname{originallemma}{Original lemma}{Original lemmas}
\makeatletter
\newlabel{lemma-sheaf-with-exact-support-quasi-coherent}{{0A78}{source}{}{sourcepointer-5}{}}
\newlabel{lemma-sheaf-with-exact-support-quasi-coherent@cref}{{[originallemma][0][]0A78}{[1][0][]source}{}{}{}}
\makeatother
\crefname{originallemma}{original lemma}{original lemmas}
\Crefname{originallemma}{Original lemma}{Original lemmas}
\makeatletter
\newlabel{lemma-shriek-closed-immersion}{{0AA2}{source}{}{sourcepointer-6}{}}
\newlabel{lemma-shriek-closed-immersion@cref}{{[originallemma][0][]0AA2}{[1][0][]source}{}{}{}}
\makeatother
\title{The upper shriek of a dualizing complex is dualizing}
\author{The Stacks Project Authors}
\date{stacks-project:duality.tex:local-snapshot-20261012; GNU FDL 1.2 or later}
\begin{document}
\maketitle
\noindent\textit{Editorial note.} All delivered mathematical text is the author's own, copied from the pinned Stacks Project chapter file duality.tex; no statement or proof step was rewritten, reordered or removed.\par\medskip
\noindent\textit{Editorial note.} Preamble provenance: the pinned chapter file opens with an input of the shared preamble and carries no preamble text of its own, so the project's single Stacks-wide preamble staged at sources/stacks-duality/preamble.tex is excerpted verbatim under the recipe fixed by LANE\_BRIEF 22.1, as three spans, lines 8-13 (verbatim package and the reference/slogan/history environments), lines 14-18 (Xy-pic setup) and lines 154-217 (the native theorem declarations and mathematical macros). Nothing was added, removed or redefined; the wrapper class and the amsmath/amssymb/amsthm/graphicx base packages are supplied by the assembler.\par\medskip
\noindent\textit{Editorial note.} The counted statement and proof contain no cross reference, no citation and no included file, so the delivered body is self contained.\par\medskip
\noindent\textit{Editorial note.} Canonical location of the counted proposition: Stacks tag 0AA3 (sources/stacks-duality/pinned-original-source.tex line 4708).\par\medskip
\noindent\textit{Allowed standard prerequisites (editorial).}\begin{itemize}
\item The shriek of the affine line and dualizingness for polynomial rings. Same chapter, Lemma labelled lemma-shriek-affine-line, cited at proof lines 4740-4742, and Dualizing Complexes, Lemma labelled dualizing-lemma-dualizing-polynomial-ring, cited at proof lines 4742-4744. Handles the projections.
\item The shriek of a closed immersion and dualizingness for quotients. Dualizing Complexes, Lemma labelled dualizing-lemma-dualizing-quotient, cited at proof lines 4746-4748, with Same chapter, Lemma labelled lemma-shriek-closed-immersion and Lemma labelled lemma-sheaf-with-exact-support-quasi-coherent, cited at proof lines 4744-4748. Handles the closed immersion.
\end{itemize}\medskip
\ifdefined\hypertarget\hypertarget{sourcepointer-1}{}\fi
\noindent\textit{Source reference.} Original reference 0F42: sources/stacks-duality/pinned-original-source.tex, source lines 3953--3953. Same-chapter pointer into the author environment situation at source line 3953; invoked at proof line None of the counted proof. It is retained only as a navigation pointer; every load-bearing step is inside the delivered span.\par\medskip
\ifdefined\hypertarget\hypertarget{sourcepointer-2}{}\fi
\noindent\textit{Source reference.} Original lemma 0AA1: sources/stacks-duality/pinned-original-source.tex, source lines 4609--4609. Same-chapter pointer into the author environment lemma at source line 4609; invoked at proof line 4721 of the counted proof. It is retained only as a navigation pointer; every load-bearing step is inside the delivered span.\par\medskip
\ifdefined\hypertarget\hypertarget{sourcepointer-3}{}\fi
\noindent\textit{Source reference.} Original lemma 0A7J: sources/stacks-chow/tags.txt, source lines 12882--12882. Cross-chapter standard prerequisite located in the project tag list staged in the local snapshot; invoked at proof line 4722 of the counted proof. It is retained only as a navigation pointer; every load-bearing step is inside the delivered span.\par\medskip
\ifdefined\hypertarget\hypertarget{sourcepointer-4}{}\fi
\noindent\textit{Source reference.} Original lemma 0A7I: sources/stacks-chow/tags.txt, source lines 12881--12881. Cross-chapter standard prerequisite located in the project tag list staged in the local snapshot; invoked at proof line 4726 of the counted proof. It is retained only as a navigation pointer; every load-bearing step is inside the delivered span.\par\medskip
\ifdefined\hypertarget\hypertarget{sourcepointer-5}{}\fi
\noindent\textit{Source reference.} Original lemma 0A78: sources/stacks-duality/pinned-original-source.tex, source lines 2373--2373. Same-chapter pointer into the author environment lemma at source line 2373; invoked at proof line 4728 of the counted proof. It is retained only as a navigation pointer; every load-bearing step is inside the delivered span.\par\medskip
\ifdefined\hypertarget\hypertarget{sourcepointer-6}{}\fi
\noindent\textit{Source reference.} Original lemma 0AA2: sources/stacks-duality/pinned-original-source.tex, source lines 4624--4624. Same-chapter pointer into the author environment lemma at source line 4624; invoked at proof line 4729 of the counted proof. It is retained only as a navigation pointer; every load-bearing step is inside the delivered span.\par\medskip
\begin{lemma}
\label{lemma-shriek-dualizing}
In Situation \ref{situation-shriek} let $f : X \to Y$ be a morphism of
$\textit{FTS}_S$. If $K$ is a dualizing complex
for $Y$, then $f^!K$ is a dualizing complex for $X$.
\end{lemma}

\begin{proof}
The question is local on $X$ hence we may assume that $X$ and $Y$ are
affine schemes. In this case we can factor $f : X \to Y$ as
$$
X \xrightarrow{i} \mathbf{A}^n_Y \to \mathbf{A}^{n - 1}_Y \to \ldots \to
\mathbf{A}^1_Y \to Y
$$
where $i$ is a closed immersion. By Lemma \ref{lemma-shriek-affine-line} and
Dualizing Complexes, Lemma \ref{dualizing-lemma-dualizing-polynomial-ring}
and induction we see that
the $p^!K$ is a dualizing complex on $\mathbf{A}^n_Y$ where
$p : \mathbf{A}^n_Y \to Y$ is the projection. Similarly, by
Dualizing Complexes, Lemma \ref{dualizing-lemma-dualizing-quotient}
and Lemmas
\ref{lemma-sheaf-with-exact-support-quasi-coherent} and
\ref{lemma-shriek-closed-immersion} we see that $i^!$
transforms dualizing complexes into dualizing complexes.
\end{proof}

\end{document}
